% BHU PYQ -- Professional Exam Template
\documentclass[11pt,a4paper]{article}

\usepackage[a4paper,top=2.5cm,bottom=2.5cm,left=2.8cm,right=2.8cm,headheight=14pt]{geometry}
\usepackage{amsmath,amssymb}
\usepackage[T1]{fontenc}
\usepackage[utf8]{inputenc}
\usepackage{lmodern}
\usepackage{microtype}
\usepackage{fancyhdr}
\usepackage{enumitem}
\usepackage{hyperref}
\usepackage{lastpage}
\usepackage{parskip}

\hypersetup{colorlinks=true,urlcolor=black,linkcolor=black,
  pdftitle={B.Sc. (Hons.) Botany Sem-III BOB-301 -- BHU PYQ}}

\pagestyle{fancy}
\fancyhf{}
\renewcommand{\headrulewidth}{0.4pt}
\renewcommand{\footrulewidth}{0.4pt}
\fancyhead[L]{\small   \textit{Banaras Hindu University}}
\fancyhead[R]{\small   \textit{Botany Paper: BOB-301}}
\fancyfoot[C]{\small Page  \thepage\ of \pageref{LastPage}}


\newlist{parts}{enumerate}{2}
\setlist[parts,1]{label=(\alph*),leftmargin=*,itemsep=4pt,topsep=3pt}
\setlist[parts,2]{label=(\roman*),leftmargin=*,itemsep=2pt,topsep=2pt}

\newcommand{\pts}[1]{\hfill   \textit{[#1]}}


\begin{document}


\begin{center}
  {\large\bfseries BANARAS HINDU UNIVERSITY}\\[2pt]
  {\normalsize B.Sc. (Hons.) Semester-III Examination, 2022-23}\\[6pt]
  \rule{0.8\linewidth}{0.5pt}\\[6pt]
  {\large\bfseries Botany}\\[4pt]
  {\large\bfseries Paper: BOB-301 --- Plant Ecology and Physiology}\\[4pt]
  {\normalsize Time: Three Hours \hfill Maximum Marks: 70}
\end{center}

\rule{\linewidth}{0.4pt}

\smallskip



\noindent   \textit{Instructions: Answer five questions in all, selecting two each from Sections A and B and Question 1 which is compulsory. The figures in the right-hand margin indicate marks.}

\bigskip

\medskip


\noindent   \textbf{Question 1.} Answer the following questions in not more than 50 words each:\pts{$1  \times 10 = 10$}

\begin{parts}
  \item What is the full form of RuBisCO?
  \item What is apical dominance?
  \item Name a pump found in the plasma membrane.
  \item Define autogenic succession.
  \item What is tropopause?
  \item Name the precursor of auxins.
  \item What are hydathodes?
  \item What is the function of cholesterol in plasma membranes?
  \item What is residual soil?
  \item What are aquaporins?
\end{parts}

\bigskip

\begin{center}
   \textbf{SECTION A}
\end{center}
\medskip

\medskip


\noindent   \textbf{Question 2.} With the help of suitable diagrams, discuss the different adaptations of Xerophytes.\pts{15}

\medskip


\noindent   \textbf{Question 3.} Write short notes on any two of the following:\pts{$7.5  \times 2 = 15$}

\begin{parts}
  \item Quantitative characters of a community
  \item Population growth-curves
  \item Ecological pyramids
\end{parts}

\medskip


\noindent   \textbf{Question 4.} Discuss in detail any two of the following:\pts{$7.5  \times 2 = 15$}

\begin{parts}
  \item Survivorship curves
  \item Soil profile
  \item Effect of temperature on the distribution of plants
\end{parts}

\medskip


\noindent   \textbf{Question 5.} Describe the different developmental stages of a hydrosere.\pts{15}

\bigskip

\begin{center}
   \textbf{SECTION B}
\end{center}
\medskip

\medskip


\noindent   \textbf{Question 6.} With the help of suitable diagram, describe the fluid mosaic model of the structure of cell membrane.\pts{15}

\medskip


\noindent   \textbf{Question 7.} Write short notes on any two of the following:\pts{$7.5  \times 2 = 15$}

\begin{parts}
  \item Components of water potential
  \item Physiological role of gibberellins
  \item Mechanism of phytochrome action
\end{parts}

\medskip


\noindent   \textbf{Question 8.} Discuss any two of the following in brief:\pts{$7.5  \times 2 = 15$}

\begin{parts}
  \item Secondary active transport
  \item Glycolysis
  \item Structure of nitrogenase
\end{parts}

\medskip


\noindent   \textbf{Question 9.} What is phosphorylation? With the help of a suitable diagram, describe the chemiosmotic model of phosphorylation in plants.\pts{15}

\vfill
\rule{\linewidth}{0.4pt}

\begin{center}\small
\href{https://bhu-exam.vercel.app/}{Click Here}\\[4pt]\href{https://me-aryan.vercel.app/}{Click Here}\\[4pt]\href{https://quashan-me.vercel.app/}{Click Here}
\end{center}

\end{document}
